\section{动态规划章正文案例推导补强}

这一节补齐本章正文出现过的 LeetCode 案例推导。阅读顺序统一按“题目说明、样例入口、暴力基线、暴力过程、重复工作、优化条件、相邻状态、状态复用、指针和字段、代码落点、正确性、边界实验、复杂度变化、迁移追问”展开；目的不是新增题量，而是把正文中的案例都讲成可复述、可证明、可迁移的解题链。

\subsection{LeetCode 70 爬楼梯}

\textbf{题目说明。} LeetCode 70 爬楼梯 的题面入口要先还原成具体任务：最后一步来自前一阶或前两阶，是最小的重叠子问题模型。

\textbf{样例入口。} 拿 n=4 手算，最后一步要么从第 3 阶跨 1 阶，要么从第 2 阶跨 2 阶，这两类互斥且覆盖所有走法。于是 ways[4]=ways[3]+ways[2]。

\textbf{暴力基线。} 暴力递归枚举每一步走 1 阶或 2 阶，直到正好到 n 计数，超过 n 则失败。

\textbf{暴力过程。} 以 n=4 为例，最后一步要么从第 3 阶走 1 阶，要么从第 2 阶走 2 阶；这两类互斥且覆盖所有走法。

\textbf{重复工作。} 题面模型：最后一步来自前一阶或前两阶，是最小的重叠子问题模型。暴力版的慢点要落到本题：浪费在递归树反复计算 ways(2)、ways(3) 等同一子问题。到达某一阶后，后面的走法数量和之前怎么来无关。后面的优化只消掉这类重复工作，不能改变答案集合。

\textbf{优化条件。} 本题的关键观察是：滚动变量保存最近两个状态即可，不需要完整数组。这句话必须能翻译成可维护状态；如果题目条件改变后观察不再成立，就不能继续套原来的写法。

\textbf{相邻状态。} 规律不是背斐波那契，而是最后一步只有两个来源，状态只需保存前两项。把这个特例继续拆开看：先问暴力搜索里哪些不同路径到达同一个后续问题，再给这个后续问题命名为状态。状态必须包含未来决策需要的全部信息，转移只从更小且已经算清的状态来。

\textbf{状态复用。} dp[i] 表示到达第 i 阶的方法数，转移是 dp[i]=dp[i-1]+dp[i-2]。再压缩成两个滚动变量。手算时只改被新输入影响的那一小部分，而不是回头重算完整候选。

\textbf{指针和字段。} 状态字段要能说清：prev2 表示 dp[i-2]，prev1 表示 dp[i-1]，current 表示 dp[i]。手算时按这个协议跟踪：n=4 时 dp[1]=1、dp[2]=2、dp[3]=3、dp[4]=5。每一步只依赖前两项。

\textbf{代码落点。} 关键是初始化：n=0 按题意通常返回 1 或不出现；LeetCode 输入 n>=1，n=1 返回 1，n=2 返回 2。关键代码行必须能逐句翻译成“旧状态怎样变成新状态”，否则就只是模板。

\textbf{正确性。} 所有到第 i 阶的方案按最后一步分成从 i-1 来和从 i-2 来两类，二者不重不漏。

\textbf{边界实验。} 先用n 为一、n 为二、题目是否允许零阶、结果范围做反例检查。实验时覆盖n 为一、n 为二、题目是否允许零阶、结果范围，先打印完整 DP 表，再比较空间压缩版本。若压缩版不同，优先检查遍历顺序和旧值保存。

\textbf{复杂度变化。} 暴力递归指数级；DP O(n) 时间，滚动变量 O(1) 空间。

\textbf{迁移追问。} 如果题目改成“若每次可走一到 k 阶，转移变成固定宽度窗口求和”，先判断原状态是否仍然覆盖新答案集合，再决定是微调边界、改 value 语义，还是换算法。

\subsection{LeetCode 198 打家劫舍}

\textbf{题目说明。} LeetCode 198 打家劫舍 的题面入口要先还原成具体任务：每间房只有偷和不偷两种结局，偷当前就不能偷前一间。

\textbf{样例入口。} 拿 [2, 7, 9] 手算，到房屋 9 时只有两类合法结局：不偷它，答案沿用前一间；偷它，就只能接前两间的最优。

\textbf{暴力基线。} 暴力递归每间房偷或不偷，偷了当前房就不能偷下一间，最后取最大金额。

\textbf{暴力过程。} 以 nums=[2, 7, 9] 为例，到房屋 9 时，合法结局只有两类：不偷 9 沿用前一间最优，偷 9 就只能接前两间最优。

\textbf{重复工作。} 题面模型：每间房只有偷和不偷两种结局，偷当前就不能偷前一间。暴力版的慢点要落到本题：浪费在递归中同一个下标之后的最优收益被反复计算。未来只取决于当前处理到哪一间。后面的优化只消掉这类重复工作，不能改变答案集合。

\textbf{优化条件。} 本题的关键观察是：滚动保存前一状态和前两状态。这句话必须能翻译成可维护状态；如果题目条件改变后观察不再成立，就不能继续套原来的写法。

\textbf{相邻状态。} 规律是 dp[i]=max(dp[i-1], dp[i-2]+nums[i])，相邻限制来自题面而不是公式本身。把这个特例继续拆开看：先问暴力搜索里哪些不同路径到达同一个后续问题，再给这个后续问题命名为状态。状态必须包含未来决策需要的全部信息，转移只从更小且已经算清的状态来。

\textbf{状态复用。} dp[i] 表示考虑前 i 间房的最大收益。转移为 max(dp[i-1], dp[i-2]+nums[i-1])，可压成两个变量。手算时只改被新输入影响的那一小部分，而不是回头重算完整候选。

\textbf{指针和字段。} 状态字段要能说清：previous\_one 是 dp[i-1]，previous\_two 是 dp[i-2]，money 是当前房屋金额。手算时按这个协议跟踪：[2, 7, 9] 中前两步最优到 7；处理 9 时比较不偷得 7、偷得 2+9=11，更新为 11。

\textbf{代码落点。} 关键是滚动更新顺序：先算 current=max(previous\_one, previous\_two+money)，再 previous\_two=previous\_one，previous\_one=current。关键代码行必须能逐句翻译成“旧状态怎样变成新状态”，否则就只是模板。

\textbf{正确性。} 最优方案对第 i 间房只有偷或不偷两类，二者互斥且覆盖全部合法方案。

\textbf{边界实验。} 先用空数组、单间、全零、金额很大做反例检查。实验时覆盖空数组、单间、全零、金额很大，先打印完整 DP 表，再比较空间压缩版本。若压缩版不同，优先检查遍历顺序和旧值保存。

\textbf{复杂度变化。} 暴力递归指数级；DP O(n) 时间，滚动 O(1) 空间。

\textbf{迁移追问。} 如果题目改成“环形房屋要拆成不含首和不含尾两个线性问题”，先判断原状态是否仍然覆盖新答案集合，再决定是微调边界、改 value 语义，还是换算法。

\subsection{LeetCode 72 编辑距离}

\textbf{题目说明。} LeetCode 72 编辑距离 的题面入口要先还原成具体任务：二维状态表示两个前缀之间的最少编辑操作。

\textbf{样例入口。} 拿 word1=ab、word2=ac 手算最后一个字符。若 b 改成 c，是 dp[1][1]+1；若删掉 b，是 dp[1][2]+1；若插入 c，是 dp[2][1]+1。字符相等时才直接继承左上。

\textbf{暴力基线。} 暴力递归比较两个字符串后缀：相等就同时跳过，不等就尝试插入、删除、替换三种操作。

\textbf{暴力过程。} 以 word1=ab、word2=ac 为例，末尾 b 和 c 不同。替换 b 成 c 来自 dp[1][1]+1；删除 b 来自 dp[1][2]+1；插入 c 来自 dp[2][1]+1。

\textbf{重复工作。} 题面模型：二维状态表示两个前缀之间的最少编辑操作。暴力版的慢点要落到本题：浪费在不同编辑序列会反复到达同一对前缀长度。只要 i 和 j 相同，后续最少编辑次数就相同。后面的优化只消掉这类重复工作，不能改变答案集合。

\textbf{优化条件。} 本题的关键观察是：末尾字符相同继承左上；不同则枚举插入、删除、替换。这句话必须能翻译成可维护状态；如果题目条件改变后观察不再成立，就不能继续套原来的写法。

\textbf{相邻状态。} 规律是每个格子枚举最后一次编辑操作，而不是凭感觉填三邻居最小。把这个特例继续拆开看：先问暴力搜索里哪些不同路径到达同一个后续问题，再给这个后续问题命名为状态。状态必须包含未来决策需要的全部信息，转移只从更小且已经算清的状态来。

\textbf{状态复用。} dp[i][j] 表示 word1 前 i 个字符变成 word2 前 j 个字符的最少操作数。按前缀长度从小到大填表。手算时只改被新输入影响的那一小部分，而不是回头重算完整候选。

\textbf{指针和字段。} 状态字段要能说清：i/j 是前缀长度；dp[i-1][j] 对应删除，dp[i][j-1] 对应插入，dp[i-1][j-1] 对应替换或字符相等继承。手算时按这个协议跟踪：空串到 ac 需要插入 2 次；ab 到空串需要删除 2 次。内部格子按末尾字符是否相等决定转移。

\textbf{代码落点。} 关键是第一行 dp[0][j]=j、第一列 dp[i][0]=i。字符相等时不要额外加一。关键代码行必须能逐句翻译成“旧状态怎样变成新状态”，否则就只是模板。

\textbf{正确性。} 任意最优编辑序列的最后一步只有插入、删除、替换或不操作四类，转移枚举了全部最后一步。

\textbf{边界实验。} 先用空串、完全相同、完全不同、操作代价不同做反例检查。实验时覆盖空串、完全相同、完全不同、操作代价不同，先打印完整 DP 表，再比较空间压缩版本。若压缩版不同，优先检查遍历顺序和旧值保存。

\textbf{复杂度变化。} 暴力递归指数级；二维 DP O(mn) 时间和空间，可滚动到 O(min(m, n)) 空间。

\textbf{迁移追问。} 如果题目改成“若只允许插入删除，问题退化到和最长公共子序列相关”，先判断原状态是否仍然覆盖新答案集合，再决定是微调边界、改 value 语义，还是换算法。

\subsection{LeetCode 312 戳气球}

\textbf{题目说明。} LeetCode 312 戳气球 的题面入口要先还原成具体任务：先戳哪个会让邻居持续变化，改成最后戳某个气球后区间边界才稳定。

\textbf{样例入口。} 拿 [3, 1, 5] 手算，若先戳 1，左右邻居会变；若改问某个区间里最后戳谁，最后一刻左右边界固定，收益可拆成左右两个子区间。

\textbf{暴力基线。} 慢版本先写递归搜索树：节点表示处理位置、剩余资源和当前选择。只要不同路径反复到达同一个节点，就说明需要把这个节点命名成 DP 状态。放到本题就是：先按“先戳哪个会让邻居持续变化，改成最后戳某个气球后区间边界才稳定”完整试慢版本；样例里已经能看到“规律是区间 DP 枚举最后戳的 mid，而不是第一个戳的气球”，所以优化前必须先能说清慢版本枚举了哪些候选。

\textbf{暴力过程。} 拿 [3, 1, 5] 手算，若先戳 1，左右邻居会变；若改问某个区间里最后戳谁，最后一刻左右边界固定，收益可拆成左右两个子区间。

\textbf{重复工作。} 题面模型：先戳哪个会让邻居持续变化，改成最后戳某个气球后区间边界才稳定。暴力版的慢点要落到本题：慢版本会在“先戳哪个会让邻居持续变化，改成最后戳某个气球后区间边界才稳定”的选择树里反复到达相同参数。样例规律是“规律是区间 DP 枚举最后戳的 mid，而不是第一个戳的气球”，说明这个参数组合可以命名成状态，算过之后后面直接复用。后面的优化只消掉这类重复工作，不能改变答案集合。

\textbf{优化条件。} 本题的关键观察是：区间 DP 定义开区间内全部戳完的最大收益，枚举最后一个被戳的气球作为切分点。这句话必须能翻译成可维护状态；如果题目条件改变后观察不再成立，就不能继续套原来的写法。

\textbf{相邻状态。} 规律是区间 DP 枚举最后戳的 mid，而不是第一个戳的气球。把这个特例继续拆开看：先问暴力搜索里哪些不同路径到达同一个后续问题，再给这个后续问题命名为状态。状态必须包含未来决策需要的全部信息，转移只从更小且已经算清的状态来。

\textbf{状态复用。} 把重复递归节点命名为状态；遍历顺序只是在保证依赖状态已经算完。本题落点是：规律是区间 DP 枚举最后戳的 mid，而不是第一个戳的气球。手算时只改被新输入影响的那一小部分，而不是回头重算完整候选。

\textbf{指针和字段。} 状态字段要能说清：状态下标、状态值语义、初始化、转移来源、遍历顺序；空间压缩时还要指出读到的是旧值还是新值。手算时按这个协议跟踪：拿 [3, 1, 5] 手算，若先戳 1，左右邻居会变；若改问某个区间里最后戳谁，最后一刻左右边界固定，收益可拆成左右两个子区间。然后按协议复查：手算时先写一个很小的 DP 表或记忆化搜索记录。每个格子旁边写它来自哪些更小状态。

\textbf{代码落点。} 先写未压缩版本校验转移；压缩前后用小表对拍；不可达哨兵参与加法前要检查溢出。关键代码行必须能逐句翻译成“旧状态怎样变成新状态”，否则就只是模板。

\textbf{正确性。} 证明从最后一步开始：任意最优解的最后一步都在转移枚举中，转移来源已经由更小状态正确计算。

\textbf{边界实验。} 先用补一的虚拟边界、区间长度遍历顺序、空区间收益为零、乘积可能较大做反例检查。实验时覆盖补一的虚拟边界、区间长度遍历顺序、空区间收益为零、乘积可能较大，先打印完整 DP 表，再比较空间压缩版本。若压缩版不同，优先检查遍历顺序和旧值保存。

\textbf{复杂度变化。} 暴力递归会重复计算相同子问题，常见指数级；DP 把每个状态算一次，时间是状态数乘单个转移成本，空间是状态表大小或压缩后的大小。

\textbf{迁移追问。} 如果题目改成“很多区间 DP 都要换成最后一步思考，避免中间操作改变边界”，先判断原状态是否仍然覆盖新答案集合，再决定是微调边界、改 value 语义，还是换算法。

\subsection{LeetCode 416 分割等和子集}

\textbf{题目说明。} LeetCode 416 分割等和子集 的题面入口要先还原成具体任务：总和一半是背包容量，每个数只能使用一次。

\textbf{样例入口。} 拿 [1, 5, 11, 5] 手算，总和 22，目标容量 11；可以用 11 单独凑成，也可以 5+5+1。

\textbf{暴力基线。} 暴力枚举每个数放到左集合还是右集合，最后检查两边和是否相等。

\textbf{暴力过程。} 以 [1, 5, 11, 5] 为例，总和 22，目标是找一个子集和为 11；可以选单独的 11，也可以选 5+5+1。

\textbf{重复工作。} 题面模型：总和一半是背包容量，每个数只能使用一次。暴力版的慢点要落到本题：浪费在不同选择顺序会反复得到相同剩余容量。问题只关心能否凑到 sum/2，不关心具体左右集合顺序。后面的优化只消掉这类重复工作，不能改变答案集合。

\textbf{优化条件。} 本题的关键观察是：0-1 背包可达性，一维容量必须倒序。这句话必须能翻译成可维护状态；如果题目条件改变后观察不再成立，就不能继续套原来的写法。

\textbf{相邻状态。} 规律是分割等和转成 0-1 背包是否能凑到 sum/2，容量必须倒序避免重复使用同一元素。把这个特例继续拆开看：阶段是物品，容量是资源，状态值是可达、最值还是计数。一个物品能否在同一轮再次使用，直接决定容量循环方向；组合和排列也由循环顺序表达。

\textbf{状态复用。} 先判断总和是否为偶数。dp[cap] 表示处理过若干数字后是否能凑出 cap，每个 num 只能用一次，所以容量倒序更新。手算时只改被新输入影响的那一小部分，而不是回头重算完整候选。

\textbf{指针和字段。} 状态字段要能说清：target=sum/2，cap 是容量，num 是当前物品；dp[cap-num] 为真则 dp[cap] 可变真。手算时按这个协议跟踪：处理 1 后可达 1；处理 5 后可达 5 和 6；处理 11 后 target 11 直接可达。

\textbf{代码落点。} 关键是 for cap from target down to num。若正序更新，会在同一轮把当前 num 用多次，变成完全背包。关键代码行必须能逐句翻译成“旧状态怎样变成新状态”，否则就只是模板。

\textbf{正确性。} 每个数只有选或不选两种状态；倒序容量保证转移读取的是上一阶段结果，不会重复使用当前数。

\textbf{边界实验。} 先用总和为奇数、目标为零、数值较大、复杂度依赖 target做反例检查。实验时覆盖总和为奇数、目标为零、数值较大、复杂度依赖 target，并故意把容量方向写反构造最小反例。这样能确认循环方向来自题意，不是背模板。

\textbf{复杂度变化。} 暴力 O(\ensuremath{2^{n}})，0-1 背包 O(n*target) 时间、O(target) 空间。

\textbf{迁移追问。} 如果题目改成“负数会破坏普通容量模型，需要重新建模”，先判断原状态是否仍然覆盖新答案集合，再决定是微调边界、改 value 语义，还是换算法。

\subsection{LeetCode 322 零钱兑换}

\textbf{题目说明。} LeetCode 322 零钱兑换 的题面入口要先还原成具体任务：金额是容量，硬币可无限使用，目标是最少硬币数。

\textbf{样例入口。} 拿 coins=[1, 2, 5]、amount=11 手算，最后一步可以用 1、2 或 5，若用 5，则来源是 dp[6]+1。

\textbf{暴力基线。} 暴力递归枚举最后使用哪枚硬币，或者枚举每种硬币使用多少次，直到金额变成 0。

\textbf{暴力过程。} 以 coins=[1, 2, 5]、amount=11 为例，若最后一枚是 5，前面必须先凑出 6；若最后一枚是 2，前面必须先凑出 9。

\textbf{重复工作。} 题面模型：金额是容量，硬币可无限使用，目标是最少硬币数。暴力版的慢点要落到本题：浪费在同一个剩余金额会从很多硬币序列到达。凑出金额 x 的最少硬币数和之前路径无关，只和 x 有关。后面的优化只消掉这类重复工作，不能改变答案集合。

\textbf{优化条件。} 本题的关键观察是：dp[value] 从 value 减 coin 的已知状态转移。这句话必须能翻译成可维护状态；如果题目条件改变后观察不再成立，就不能继续套原来的写法。

\textbf{相邻状态。} 规律是 dp[x] 表示凑成 x 的最少硬币数，完全背包允许同一硬币重复使用；不可达状态不能参与加一。把这个特例继续拆开看：阶段是物品，容量是资源，状态值是可达、最值还是计数。一个物品能否在同一轮再次使用，直接决定容量循环方向；组合和排列也由循环顺序表达。

\textbf{状态复用。} dp[x] 表示凑出金额 x 的最少硬币数，初始化 dp[0]=0、其他为 INF。对每个金额尝试所有 coin。手算时只改被新输入影响的那一小部分，而不是回头重算完整候选。

\textbf{指针和字段。} 状态字段要能说清：x 是当前金额，coin 是最后一枚硬币，dp[x-coin]+1 是使用 coin 的候选答案。手算时按这个协议跟踪：dp[5] 可由 coin=5 得到 1；dp[11] 可由 dp[6]+1、dp[9]+1、dp[10]+1 取最小。

\textbf{代码落点。} 关键是不可达 INF 不能直接加一造成溢出；完全背包一维写法可金额正序，也可金额外层逐项求最小。关键代码行必须能逐句翻译成“旧状态怎样变成新状态”，否则就只是模板。

\textbf{正确性。} 任意最优方案都有一枚最后硬币 coin，去掉它后剩余金额 x-coin 必须也是最优子问题，否则可替换变更优。

\textbf{边界实验。} 先用无法凑出、amount 为零、硬币面额大于目标、哨兵溢出做反例检查。实验时覆盖无法凑出、amount 为零、硬币面额大于目标、哨兵溢出，并故意把容量方向写反构造最小反例。这样能确认循环方向来自题意，不是背模板。

\textbf{复杂度变化。} 暴力指数级；DP O(amount*coins.size()) 时间，O(amount) 空间。

\textbf{迁移追问。} 如果题目改成“零钱兑换 II 问组合数，循环语义不同”，先判断原状态是否仍然覆盖新答案集合，再决定是微调边界、改 value 语义，还是换算法。

\subsection{LeetCode 300 最长递增子序列}

\textbf{题目说明。} LeetCode 300 最长递增子序列 的题面入口要先还原成具体任务：二次 DP 固定结尾，二分优化维护每个长度的最小结尾。

\textbf{样例入口。} 拿 [10, 9, 2, 5, 3] 手算。二次 DP 固定每个位置为结尾，5 可接在 2 后；优化时 tails[长度] 保存该长度递增子序列的最小尾值，尾值越小未来越容易接。

\textbf{暴力基线。} 暴力递归或 DP 枚举每个元素作为子序列结尾，向前寻找所有更小元素。

\textbf{暴力过程。} 以 [10, 9, 2, 5, 3, 7, 101, 18] 为例，二次 DP 可以算出以 7 结尾的长度来自 2、5 或 3 后面接 7。

\textbf{重复工作。} 题面模型：二次 DP 固定结尾，二分优化维护每个长度的最小结尾。暴力版的慢点要落到本题：二次 DP 慢在每个位置都回看所有历史元素。更快写法不需要保存真实序列，只保存每个长度的最小可能结尾。后面的优化只消掉这类重复工作，不能改变答案集合。

\textbf{优化条件。} 本题的关键观察是：tails 不是实际答案序列，而是未来扩展潜力表。这句话必须能翻译成可维护状态；如果题目条件改变后观察不再成立，就不能继续套原来的写法。

\textbf{相邻状态。} 规律是二分替换 tails，不代表真实序列完整内容，而是维护未来潜力。把这个特例继续拆开看：先问暴力搜索里哪些不同路径到达同一个后续问题，再给这个后续问题命名为状态。状态必须包含未来决策需要的全部信息，转移只从更小且已经算清的状态来。

\textbf{状态复用。} tails[len] 表示长度为 len+1 的递增子序列中最小结尾。对每个 x，用 lower\_bound 找第一个 >=x 的位置并替换。手算时只改被新输入影响的那一小部分，而不是回头重算完整候选。

\textbf{指针和字段。} 状态字段要能说清：tails 不是最终答案序列，而是未来扩展潜力表；tails.size() 是当前最长长度。手算时按这个协议跟踪：处理 2、5、3 时，tails 从 [2, 5] 变成 [2, 3]，长度没变，但结尾更小，未来更容易接 7。

\textbf{代码落点。} 严格递增用 lower\_bound；若题目改成非递减，要改用 upper\_bound。要还原路径需要额外 predecessor 数组。关键代码行必须能逐句翻译成“旧状态怎样变成新状态”，否则就只是模板。

\textbf{正确性。} 同样长度下，结尾越小，未来能接上的数字集合只会更多；替换 tails 不改变已知长度，只提升扩展潜力。

\textbf{边界实验。} 先用重复元素、严格递增和非递减区别、空数组做反例检查。实验时覆盖重复元素、严格递增和非递减区别、空数组，先打印完整 DP 表，再比较空间压缩版本。若压缩版不同，优先检查遍历顺序和旧值保存。

\textbf{复杂度变化。} 二次 DP O(\ensuremath{n^{2}})，tails 加二分 O(n log n) 时间、O(n) 空间。

\textbf{迁移追问。} 如果题目改成“若要还原路径，需要额外前驱数组，不能只看 tails”，先判断原状态是否仍然覆盖新答案集合，再决定是微调边界、改 value 语义，还是换算法。

\subsection{LeetCode 139 单词拆分}

\textbf{题目说明。} LeetCode 139 单词拆分 的题面入口要先还原成具体任务：状态表示前缀能否由字典单词拼出。

\textbf{样例入口。} 拿 s=leetcode、字典 [leet, code] 手算，dp[4] 因为 leet 可拆而为真，dp[8] 又因为 dp[4] 且 code 在字典里为真。

\textbf{暴力基线。} 暴力递归从字符串开头切一个字典单词，剩余后缀继续递归；尝试所有切法。

\textbf{暴力过程。} 以 leetcode 和字典 [leet, code] 为例，前缀 leet 在字典中，剩余 code 也在字典中，所以可拆。

\textbf{重复工作。} 题面模型：状态表示前缀能否由字典单词拼出。暴力版的慢点要落到本题：浪费在同一个后缀会由不同切法反复判断。例如某个下标 i 之后能否拆分，只取决于 i，不取决于前面怎么切。后面的优化只消掉这类重复工作，不能改变答案集合。

\textbf{优化条件。} 本题的关键观察是：枚举最后一个单词的起点，左侧前缀可达且子串在字典中则当前可达。这句话必须能翻译成可维护状态；如果题目条件改变后观察不再成立，就不能继续套原来的写法。

\textbf{相邻状态。} 规律是 dp[i] 表示前 i 个字符能否拆分，转移枚举最后一个单词起点。把这个特例继续拆开看：先问暴力搜索里哪些不同路径到达同一个后续问题，再给这个后续问题命名为状态。状态必须包含未来决策需要的全部信息，转移只从更小且已经算清的状态来。

\textbf{状态复用。} dp[i] 表示 s[0..i) 是否能被字典拆分。若存在 j<i，使 dp[j] 为真且 s[j..i) 在字典中，则 dp[i]=true。手算时只改被新输入影响的那一小部分，而不是回头重算完整候选。

\textbf{指针和字段。} 状态字段要能说清：i 是当前前缀右端，j 是最后一个单词起点，word\_set 用于 O(1) 查询子串是否为单词。手算时按这个协议跟踪：leetcode 中 dp[4] 因 leet 为真；随后 i=8 时 j=4，dp[4] 为真且 code 在字典中，所以 dp[8]=true。

\textbf{代码落点。} 关键是 dp[0]=true 表示空前缀可拆。热循环里 substr 会复制字符串，工程版可按最大单词长度限制 j 或用 Trie。关键代码行必须能逐句翻译成“旧状态怎样变成新状态”，否则就只是模板。

\textbf{正确性。} 任意合法拆分都有最后一个单词 s[j..i)，其左侧前缀必须可拆；转移枚举了所有可能最后单词起点。

\textbf{边界实验。} 先用空串、重复单词、长单词、substr 复制成本做反例检查。实验时覆盖空串、重复单词、长单词、substr 复制成本，先打印完整 DP 表，再比较空间压缩版本。若压缩版不同，优先检查遍历顺序和旧值保存。

\textbf{复杂度变化。} 暴力指数级；DP 有 O(\ensuremath{n^{2}}) 个切分检查，若 substr 计入成本可到 O(\ensuremath{n^{3}})，空间 O(n+字典)。

\textbf{迁移追问。} 如果题目改成“若要输出所有句子，需要 DFS 加记忆化或前驱列表”，先判断原状态是否仍然覆盖新答案集合，再决定是微调边界、改 value 语义，还是换算法。

\subsection{LeetCode 1143 最长公共子序列}

\textbf{题目说明。} LeetCode 1143 最长公共子序列 的题面入口要先还原成具体任务：状态表示两个前缀的最长公共子序列长度。

\textbf{样例入口。} 拿 abcde 和 ace 手算，末尾 e 相等时答案来自两个前缀 abcd 和 ac 加一；若末尾不同，就取删掉一侧末尾后的较大值。

\textbf{暴力基线。} 暴力递归枚举两个字符串的子序列选择，或从末尾开始相等就选、不等就分支删一边。

\textbf{暴力过程。} 以 text1=abcde、text2=ace 为例，末尾 e 相等时，答案来自 abcd 和 ac 的 LCS 加一；若末尾不同，就只能丢掉一侧末尾试更优。

\textbf{重复工作。} 题面模型：状态表示两个前缀的最长公共子序列长度。暴力版的慢点要落到本题：浪费在很多前缀对会被重复计算。只要两个前缀长度 i、j 固定，它们的 LCS 长度就是固定状态。后面的优化只消掉这类重复工作，不能改变答案集合。

\textbf{优化条件。} 本题的关键观察是：末尾相同取左上加一，不同取上方和左方最大。这句话必须能翻译成可维护状态；如果题目条件改变后观察不再成立，就不能继续套原来的写法。

\textbf{相邻状态。} 规律是 dp[i][j] 表示两个前缀的 LCS 长度，转移只看左、上、左上。把这个特例继续拆开看：先问暴力搜索里哪些不同路径到达同一个后续问题，再给这个后续问题命名为状态。状态必须包含未来决策需要的全部信息，转移只从更小且已经算清的状态来。

\textbf{状态复用。} dp[i][j] 表示 text1 前 i 个字符和 text2 前 j 个字符的 LCS 长度。相等看左上加一，不等取上方和左方最大。手算时只改被新输入影响的那一小部分，而不是回头重算完整候选。

\textbf{指针和字段。} 状态字段要能说清：i/j 是前缀长度，dp[i-1][j] 表示丢 text1 末尾，dp[i][j-1] 表示丢 text2 末尾，dp[i-1][j-1] 表示两个末尾配对。手算时按这个协议跟踪：a 与 a 相等得到 1；到 abcde 和 ace 的最后，e 与 e 相等，所以由 abcd/ac 的答案加一。

\textbf{代码落点。} 关键是表多开一行一列表示空前缀，循环从 1 开始，访问字符用 text[i-1]。关键代码行必须能逐句翻译成“旧状态怎样变成新状态”，否则就只是模板。

\textbf{正确性。} 任意 LCS 对末尾字符只有配对或至少丢掉一边两类；转移覆盖这两类并取最大。

\textbf{边界实验。} 先用空串、完全相同、无公共字符、子序列不是子串做反例检查。实验时覆盖空串、完全相同、无公共字符、子序列不是子串，先打印完整 DP 表，再比较空间压缩版本。若压缩版不同，优先检查遍历顺序和旧值保存。

\textbf{复杂度变化。} 暴力指数级；DP O(mn) 时间和空间，可滚动到 O(min(m, n)) 空间。

\textbf{迁移追问。} 如果题目改成“若要求还原序列，需要从表右下角反向追踪选择”，先判断原状态是否仍然覆盖新答案集合，再决定是微调边界、改 value 语义，还是换算法。

\subsection{LeetCode 518 零钱兑换 II}

\textbf{题目说明。} LeetCode 518 零钱兑换 II 的题面入口要先还原成具体任务：硬币无限使用，问组合数而不是最少数量。

\textbf{样例入口。} 拿 amount=5、coins=[1, 2, 5] 手算，组合 5、2+2+1、2+1+1+1、全 1 一共 4 种；1+2 和 2+1 不能重复算。

\textbf{暴力基线。} 暴力 DFS 枚举每种硬币使用几个，或者枚举所有硬币排列再去重。

\textbf{暴力过程。} amount=5、coins=[1, 2, 5] 时，组合有 5、2+2+1、2+1+1+1、1+1+1+1+1，共 4 种；1+2 和 2+1 不能重复算。

\textbf{重复工作。} 题面模型：硬币无限使用，问组合数而不是最少数量。暴力版的慢点要落到本题：浪费在排列顺序重复。组合计数应固定硬币顺序，让每种组合只在一种顺序下被统计。后面的优化只消掉这类重复工作，不能改变答案集合。

\textbf{优化条件。} 本题的关键观察是：外层遍历硬币、内层金额正序，保证组合按固定硬币顺序统计。这句话必须能翻译成可维护状态；如果题目条件改变后观察不再成立，就不能继续套原来的写法。

\textbf{相邻状态。} 规律是组合计数时外层枚举硬币，内层容量正序。把这个特例继续拆开看：阶段是物品，容量是资源，状态值是可达、最值还是计数。一个物品能否在同一轮再次使用，直接决定容量循环方向；组合和排列也由循环顺序表达。

\textbf{状态复用。} 完全背包计数。dp[amount] 表示当前已处理硬币集合能凑出该金额的组合数；外层枚举 coin，内层金额正序。手算时只改被新输入影响的那一小部分，而不是回头重算完整候选。

\textbf{指针和字段。} 状态字段要能说清：coin 是当前允许使用的硬币，cap 是当前金额，dp[cap-coin] 是加入一枚 coin 前的组合数。手算时按这个协议跟踪：处理 coin=2 时，dp[4] 会从 dp[2] 来，表示使用两个 2；不会再把 1+2 和 2+1 分开统计。

\textbf{代码落点。} 关键循环是 for coin in coins: for cap from coin to amount: dp[cap]+=dp[cap-coin]。若外层金额内层硬币，会统计排列数。关键代码行必须能逐句翻译成“旧状态怎样变成新状态”，否则就只是模板。

\textbf{正确性。} 外层固定硬币顺序后，每个组合按最大硬币处理阶段唯一计入；金额正序允许当前硬币重复使用。

\textbf{边界实验。} 先用amount 为零、无硬币、组合数较大、循环顺序反例做反例检查。实验时覆盖amount 为零、无硬币、组合数较大、循环顺序反例，并故意把容量方向写反构造最小反例。这样能确认循环方向来自题意，不是背模板。

\textbf{复杂度变化。} 暴力指数级；完全背包计数 O(amount*coins.size()) 时间，O(amount) 空间。

\textbf{迁移追问。} 如果题目改成“若外层金额内层硬币，会统计排列数”，先判断原状态是否仍然覆盖新答案集合，再决定是微调边界、改 value 语义，还是换算法。

\subsection{LeetCode 494 目标和}

\textbf{题目说明。} LeetCode 494 目标和 的题面入口要先还原成具体任务：正负号选择可代数转换为选若干数凑出特定正和。

\textbf{样例入口。} 拿 [1, 1, 1, 1, 1]、target=3 手算，给其中一个 1 加负号，其余为正可以得到 3，共有 5 种位置选择。也可转成选正数子集和为 (sum+target)/2。

\textbf{暴力基线。} 暴力递归给每个数添加 + 或 -，枚举 \ensuremath{2^{n}} 个符号序列并统计和为 target 的方案。

\textbf{暴力过程。} 以 [1, 1, 1, 1, 1]、target=3 为例，只有一个 1 取负、其余取正时可得到 3，共 5 种位置选择。

\textbf{重复工作。} 题面模型：正负号选择可代数转换为选若干数凑出特定正和。暴力版的慢点要落到本题：浪费在符号序列会反复得到相同部分和。把正数集合 P、负数集合 N 写成 sum(P)-sum(N)=target，可转成子集和计数。后面的优化只消掉这类重复工作，不能改变答案集合。

\textbf{优化条件。} 本题的关键观察是：若 S 加 target 为奇数或目标绝对值过大，直接无解。这句话必须能翻译成可维护状态；如果题目条件改变后观察不再成立，就不能继续套原来的写法。

\textbf{相邻状态。} 规律是符号分配变成 0-1 计数背包，奇偶和范围先判定。把这个特例继续拆开看：阶段是物品，容量是资源，状态值是可达、最值还是计数。一个物品能否在同一轮再次使用，直接决定容量循环方向；组合和排列也由循环顺序表达。

\textbf{状态复用。} 总和 total。若 total+target 为奇数或 abs(target)>total，直接 0；否则找子集和 positive=(total+target)/2 的方案数。手算时只改被新输入影响的那一小部分，而不是回头重算完整候选。

\textbf{指针和字段。} 状态字段要能说清：dp[cap] 是凑出 cap 的方案数，num 是当前数；0-1 计数背包容量倒序更新。手算时按这个协议跟踪：target=3、total=5，positive=4；选择四个 1 作为正数等价于选择哪个 1 不放进去，共 5 种。

\textbf{代码落点。} 关键是 dp[0]=1，for cap from positive down to num。数字 0 会让方案数翻倍，倒序更新能自然处理。关键代码行必须能逐句翻译成“旧状态怎样变成新状态”，否则就只是模板。

\textbf{正确性。} 每个符号分配和一个正数子集一一对应；子集和满足 positive 时，正负差正好为 target。

\textbf{边界实验。} 先用零元素会让方案数翻倍、目标为负、容量为零做反例检查。实验时覆盖零元素会让方案数翻倍、目标为负、容量为零，并故意把容量方向写反构造最小反例。这样能确认循环方向来自题意，不是背模板。

\textbf{复杂度变化。} 暴力 O(\ensuremath{2^{n}})，背包 O(n*positive) 时间，O(positive) 空间。

\textbf{迁移追问。} 如果题目改成“也可用哈希 DP 保存所有可能和，但背包转换更高效”，先判断原状态是否仍然覆盖新答案集合，再决定是微调边界、改 value 语义，还是换算法。

\subsection{LeetCode 714 买卖股票含手续费}

\textbf{题目说明。} LeetCode 714 买卖股票含手续费 的题面入口要先还原成具体任务：手续费改变交易收益，仍可用持有和不持有状态。

\textbf{样例入口。} 拿 [1, 3, 2, 8, 4, 9]、fee=2 手算，8 卖出时净赚 5；手续费让频繁小波动不一定值得交易。

\textbf{暴力基线。} 慢版本先写递归搜索树：节点表示处理位置、剩余资源和当前选择。只要不同路径反复到达同一个节点，就说明需要把这个节点命名成 DP 状态。放到本题就是：先按“手续费改变交易收益，仍可用持有和不持有状态”完整试慢版本；样例里已经能看到“规律是每天维护 hold 和 cash，卖出时扣手续费，买入时从 cash 转到 hold”，所以优化前必须先能说清慢版本枚举了哪些候选。

\textbf{暴力过程。} 拿 [1, 3, 2, 8, 4, 9]、fee=2 手算，8 卖出时净赚 5；手续费让频繁小波动不一定值得交易。

\textbf{重复工作。} 题面模型：手续费改变交易收益，仍可用持有和不持有状态。暴力版的慢点要落到本题：慢版本会在“手续费改变交易收益，仍可用持有和不持有状态”的选择树里反复到达相同参数。样例规律是“规律是每天维护 hold 和 cash，卖出时扣手续费，买入时从 cash 转到 hold”，说明这个参数组合可以命名成状态，算过之后后面直接复用。后面的优化只消掉这类重复工作，不能改变答案集合。

\textbf{优化条件。} 本题的关键观察是：卖出时扣费或买入时扣费都可以，但不要重复扣。这句话必须能翻译成可维护状态；如果题目条件改变后观察不再成立，就不能继续套原来的写法。

\textbf{相邻状态。} 规律是每天维护 hold 和 cash，卖出时扣手续费，买入时从 cash 转到 hold。把这个特例继续拆开看：先问暴力搜索里哪些不同路径到达同一个后续问题，再给这个后续问题命名为状态。状态必须包含未来决策需要的全部信息，转移只从更小且已经算清的状态来。

\textbf{状态复用。} 把重复递归节点命名为状态；遍历顺序只是在保证依赖状态已经算完。本题落点是：规律是每天维护 hold 和 cash，卖出时扣手续费，买入时从 cash 转到 hold。手算时只改被新输入影响的那一小部分，而不是回头重算完整候选。

\textbf{指针和字段。} 状态字段要能说清：状态下标、状态值语义、初始化、转移来源、遍历顺序；空间压缩时还要指出读到的是旧值还是新值。手算时按这个协议跟踪：拿 [1, 3, 2, 8, 4, 9]、fee=2 手算，8 卖出时净赚 5；手续费让频繁小波动不一定值得交易。然后按协议复查：手算时先写一个很小的 DP 表或记忆化搜索记录。每个格子旁边写它来自哪些更小状态。

\textbf{代码落点。} 先写未压缩版本校验转移；压缩前后用小表对拍；不可达哨兵参与加法前要检查溢出。关键代码行必须能逐句翻译成“旧状态怎样变成新状态”，否则就只是模板。

\textbf{正确性。} 证明从最后一步开始：任意最优解的最后一步都在转移枚举中，转移来源已经由更小状态正确计算。

\textbf{边界实验。} 先用手续费为零、价格震荡、小利润不足手续费做反例检查。实验时覆盖手续费为零、价格震荡、小利润不足手续费，先打印完整 DP 表，再比较空间压缩版本。若压缩版不同，优先检查遍历顺序和旧值保存。

\textbf{复杂度变化。} 暴力递归会重复计算相同子问题，常见指数级；DP 把每个状态算一次，时间是状态数乘单个转移成本，空间是状态表大小或压缩后的大小。

\textbf{迁移追问。} 如果题目改成“这是状态机比局部贪心更稳的例子”，先判断原状态是否仍然覆盖新答案集合，再决定是微调边界、改 value 语义，还是换算法。

\subsection{LeetCode 847 访问所有节点的最短路径}

\textbf{题目说明。} LeetCode 847 访问所有节点的最短路径 的题面入口要先还原成具体任务：访问集合影响未来能力，同一节点配不同 mask 不是同一状态。

\textbf{样例入口。} 拿 graph=[[1, 2, 3], [0], [0], [0]] 手算，从中心 0 出发访问三个叶子需要来回走，长度为 4；如果只按节点 visited，到 0 时会把访问了不同叶子的状态混掉。

\textbf{暴力基线。} 慢版本先按题意画出完整状态图，用普通搜索跑通可达性。这个过程用于确认访问集合影响未来能力，同一节点配不同 mask 不是同一状态的节点和边，之后再讨论访问标记、层次或代价顺序。放到本题就是：先按“访问集合影响未来能力，同一节点配不同 mask 不是同一状态”完整试慢版本；样例里已经能看到“规律是 BFS 状态必须是 (node, mask)，从所有起点多源出发，第一次到全集 mask 就是最短”，所以优化前必须先能说清慢版本枚举了哪些候选。

\textbf{暴力过程。} 拿 graph=[[1, 2, 3], [0], [0], [0]] 手算，从中心 0 出发访问三个叶子需要来回走，长度为 4；如果只按节点 visited，到 0 时会把访问了不同叶子的状态混掉。

\textbf{重复工作。} 题面模型：访问集合影响未来能力，同一节点配不同 mask 不是同一状态。暴力版的慢点要落到本题：慢版本会围绕“访问集合影响未来能力，同一节点配不同 mask 不是同一状态”反复展开同一批状态。样例规律是“规律是 BFS 状态必须是 (node, mask)，从所有起点多源出发，第一次到全集 mask 就是最短”，说明必须把节点、边、访问标记和资源维度一次建清；同一状态不应重复入队，不同资源状态也不能误合并。后面的优化只消掉这类重复工作，不能改变答案集合。

\textbf{优化条件。} 本题的关键观察是：从所有节点以各自单点 mask 多源入队，BFS 扩展 (node, mask)，第一次到全集 mask 即最短。这句话必须能翻译成可维护状态；如果题目条件改变后观察不再成立，就不能继续套原来的写法。

\textbf{相邻状态。} 规律是 BFS 状态必须是 (node, mask)，从所有起点多源出发，第一次到全集 mask 就是最短。把这个特例继续拆开看：状态节点不一定等于原始位置，还可能包含钥匙、剩余资源、时间或访问集合。visited 只能合并未来能力完全相同的状态，否则会把合法路径剪掉。

\textbf{状态复用。} 把重复搜索压成访问标记、距离数组或拓扑状态；邻居生成也要从全量扫描变成结构化枚举。本题落点是：规律是 BFS 状态必须是 (node, mask)，从所有起点多源出发，第一次到全集 mask 就是最短。手算时只改被新输入影响的那一小部分，而不是回头重算完整候选。

\textbf{指针和字段。} 状态字段要能说清：状态编号、邻接生成方式、访问标记、距离或层数、队列内容；若状态含资源，visited 不能只按位置记录。手算时按这个协议跟踪：拿 graph=[[1, 2, 3], [0], [0], [0]] 手算，从中心 0 出发访问三个叶子需要来回走，长度为 4；如果只按节点 visited，到 0 时会把访问了不同叶子的状态混掉。然后按协议复查：手算时画队列或栈，状态必须包含题目需要的所有资源。入队时标记还是出队时标记，要和去重语义一致。

\textbf{代码落点。} 入队时标记还是出队时标记必须一致；方向数组集中定义；多源 BFS 要一次性把所有源点放入初始队列。关键代码行必须能逐句翻译成“旧状态怎样变成新状态”，否则就只是模板。

\textbf{正确性。} 证明从可达性开始：搜索覆盖所有合法边能到达的状态，访问标记只删除重复状态而不删除不同未来能力。

\textbf{边界实验。} 先用单节点、完全图、星形图、重复经过节点、visited 必须包含 mask做反例检查。实验时覆盖单节点、完全图、星形图、重复经过节点、visited 必须包含 mask，并打印入队时的完整状态。若同一位置但资源不同的状态被合并，很多图题会出现隐蔽错误。

\textbf{复杂度变化。} 暴力重复从多个起点搜索会反复遍历图；正确建模和访问标记后，BFS 或 DFS 通常按节点数加边数计算，状态图题则按完整状态数计算。

\textbf{迁移追问。} 如果题目改成“若边有权，状态图要改用 Dijkstra”，先判断原状态是否仍然覆盖新答案集合，再决定是微调边界、改 value 语义，还是换算法。

\subsection{LeetCode 53 最大子数组和}

\textbf{题目说明。} LeetCode 53 最大子数组和给定整数数组 nums，要求找一个非空连续子数组，使元素和最大，并返回这个最大和。

\textbf{样例入口。} 样例 nums=[-2, 1, -3, 4, -1, 2, 1, -5, 4]，输出 6；连续段 [4, -1, 2, 1] 的和为 6。

\textbf{暴力基线。} 暴力枚举所有连续子数组，累计每段和并更新最大值；优化一点可以固定左端、右端扩展时维护当前和。

\textbf{暴力过程。} 以 [-2, 1, -3, 4, -1, 2, 1] 为例，暴力会重复计算很多以 4 开头或以 2 结尾的子段。真正关键是当前数要不要接上前一个后缀。

\textbf{重复工作。} 题面模型：给定整数数组，返回一个非空连续子数组能够取得的最大元素和。暴力版的慢点要落到本题：浪费在重复求每个后缀和。对于以 i 结尾的最优连续段，历史信息只需要一个值：以 i-1 结尾的最大后缀和。后面的优化只消掉这类重复工作，不能改变答案集合。

\textbf{优化条件。} 本题的关键观察是：用 [-2, 3] 和 [2, 3] 对比，推出当前位置只在单独开始与接上旧后缀之间取较大值。这句话必须能翻译成可维护状态；如果题目条件改变后观察不再成立，就不能继续套原来的写法。

\textbf{相邻状态。} 规律是当前位置只需要比较 nums[i] 与 bestEndingHere + nums[i]，全局答案再和 bestEndingHere 比较；全负数组要求答案初始化为第一个数，不能初始化为 0。把这个特例继续拆开看：先问暴力搜索里哪些不同路径到达同一个后续问题，再给这个后续问题命名为状态。状态必须包含未来决策需要的全部信息，转移只从更小且已经算清的状态来。

\textbf{状态复用。} bestEndingHere 表示以当前位置结尾的最大子数组和；answer 表示全局最大值。当前位置要么单独开始，要么接上旧后缀。手算时只改被新输入影响的那一小部分，而不是回头重算完整候选。

\textbf{指针和字段。} 状态字段要能说清：current=bestEndingHere，answer 是全局最优；nums[i] 与 current+nums[i] 比较决定是否丢弃旧后缀。手算时按这个协议跟踪：遇到 4 之前，旧后缀为负，接上反而变差，所以从 4 重新开始；之后 -1、2、1 继续接上得到 6。

\textbf{代码落点。} 关键代码是 current=max(nums[i], current+nums[i]); answer=max(answer, current)。answer 必须初始化为 nums[0]，否则全负数组会错。关键代码行必须能逐句翻译成“旧状态怎样变成新状态”，否则就只是模板。

\textbf{正确性。} 任意以 i 结尾的连续段只有两类：只含 nums[i]，或由某个以 i-1 结尾的连续段加上 nums[i]。取前一类最优即可覆盖全部候选。

\textbf{边界实验。} 先用全负数组、单元素、和溢出、答案不能初始化为零做反例检查。实验时覆盖全负数组、单元素、和溢出、答案不能初始化为零，先打印完整 DP 表，再比较空间压缩版本。若压缩版不同，优先检查遍历顺序和旧值保存。

\textbf{复杂度变化。} 暴力 O(\ensuremath{n^{2}}) 或 O(\ensuremath{n^{3}})，Kadane 算法 O(n) 时间、O(1) 空间。

\textbf{迁移追问。} 如果题目改成“若要求返回区间下标，同步记录重新开始位置和最优左右端”，先判断原状态是否仍然覆盖新答案集合，再决定是微调边界、改 value 语义，还是换算法。

\subsection{LeetCode 377 组合总和 IV}

\textbf{题目说明。} LeetCode 377 组合总和 IV 的题面入口要先还原成具体任务：题目统计排列数，选择顺序不同算不同方案。

\textbf{样例入口。} 拿 nums=[1, 2, 3]、target=4 手算，1+3 和 3+1 是不同排列，所以 dp[4] 要按最后一步来自 1、2、3 分别累加。

\textbf{暴力基线。} 慢版本枚举每个物品选不选、选几次，并记录阶段和剩余容量。相同阶段和容量反复出现时，就可以把指数搜索压成表。放到本题就是：先按“题目统计排列数，选择顺序不同算不同方案”完整试慢版本；样例里已经能看到“规律是组合总和 IV 计数排列，容量在外层、数字在内层，循环顺序不能写成组合题”，所以优化前必须先能说清慢版本枚举了哪些候选。

\textbf{暴力过程。} 拿 nums=[1, 2, 3]、target=4 手算，1+3 和 3+1 是不同排列，所以 dp[4] 要按最后一步来自 1、2、3 分别累加。

\textbf{重复工作。} 题面模型：题目统计排列数，选择顺序不同算不同方案。暴力版的慢点要落到本题：慢版本会围绕“题目统计排列数，选择顺序不同算不同方案”枚举物品使用情况。样例规律是“规律是组合总和 IV 计数排列，容量在外层、数字在内层，循环顺序不能写成组合题”，说明阶段、容量和使用次数已经足够描述后续选择，没必要反复展开同一剩余容量。后面的优化只消掉这类重复工作，不能改变答案集合。

\textbf{优化条件。} 本题的关键观察是：外层遍历容量，内层遍历数字，让每个容量的方案按最后一个数累加。这句话必须能翻译成可维护状态；如果题目条件改变后观察不再成立，就不能继续套原来的写法。

\textbf{相邻状态。} 规律是组合总和 IV 计数排列，容量在外层、数字在内层，循环顺序不能写成组合题。把这个特例继续拆开看：阶段是物品，容量是资源，状态值是可达、最值还是计数。一个物品能否在同一轮再次使用，直接决定容量循环方向；组合和排列也由循环顺序表达。

\textbf{状态复用。} 把选择树压成阶段和容量；物品是否可重复使用决定一维表的遍历方向。本题落点是：规律是组合总和 IV 计数排列，容量在外层、数字在内层，循环顺序不能写成组合题。手算时只改被新输入影响的那一小部分，而不是回头重算完整候选。

\textbf{指针和字段。} 状态字段要能说清：物品阶段、容量、状态值、循环方向、取模或不可达哨兵；组合和排列必须用不同循环顺序表达。手算时按这个协议跟踪：拿 nums=[1, 2, 3]、target=4 手算，1+3 和 3+1 是不同排列，所以 dp[4] 要按最后一步来自 1、2、3 分别累加。然后按协议复查：手算时用一个容量很小的例子比较正序和倒序。只要循环方向改变后同一物品能否重复使用发生变化，就说明方向不是细节。

\textbf{代码落点。} 0-1 倒序、完全正序、排列外层容量、组合外层物品；每种循环方向都要能用一个小例子解释。关键代码行必须能逐句翻译成“旧状态怎样变成新状态”，否则就只是模板。

\textbf{正确性。} 证明从当前物品使用次数开始：0-1、完全、计数和最值的循环语义必须和题意一致。

\textbf{边界实验。} 先用target 为零、数字大于目标、方案数溢出、循环顺序做反例检查。实验时覆盖target 为零、数字大于目标、方案数溢出、循环顺序，并故意把容量方向写反构造最小反例。这样能确认循环方向来自题意，不是背模板。

\textbf{复杂度变化。} 暴力枚举物品选择指数级；背包 DP 通常按物品数乘容量或多维容量乘积计算，这是伪多项式复杂度。

\textbf{迁移追问。} 如果题目改成“若顺序不同不算新方案，就变成完全背包组合数，循环顺序要交换”，先判断原状态是否仍然覆盖新答案集合，再决定是微调边界、改 value 语义，还是换算法。

\subsection{LeetCode 309 最佳买卖股票含冷冻期}

\textbf{题目说明。} LeetCode 309 最佳买卖股票含冷冻期 的题面入口要先还原成具体任务：冷冻期让买入来源受昨天卖出状态限制。

\textbf{样例入口。} 拿 [1, 2, 3, 0, 2] 手算，第 2 天卖出后第 3 天不能立刻买入，因为有冷冻期；状态必须区分持股、刚卖出、空仓可买。

\textbf{暴力基线。} 慢版本先写递归搜索树：节点表示处理位置、剩余资源和当前选择。只要不同路径反复到达同一个节点，就说明需要把这个节点命名成 DP 状态。放到本题就是：先按“冷冻期让买入来源受昨天卖出状态限制”完整试慢版本；样例里已经能看到“规律是冷冻期把普通买卖股票拆成多个日末状态，转移不能跨过限制”，所以优化前必须先能说清慢版本枚举了哪些候选。

\textbf{暴力过程。} 拿 [1, 2, 3, 0, 2] 手算，第 2 天卖出后第 3 天不能立刻买入，因为有冷冻期；状态必须区分持股、刚卖出、空仓可买。

\textbf{重复工作。} 题面模型：冷冻期让买入来源受昨天卖出状态限制。暴力版的慢点要落到本题：慢版本会在“冷冻期让买入来源受昨天卖出状态限制”的选择树里反复到达相同参数。样例规律是“规律是冷冻期把普通买卖股票拆成多个日末状态，转移不能跨过限制”，说明这个参数组合可以命名成状态，算过之后后面直接复用。后面的优化只消掉这类重复工作，不能改变答案集合。

\textbf{优化条件。} 本题的关键观察是：状态机维护持有、刚卖出、休息三种日终状态。这句话必须能翻译成可维护状态；如果题目条件改变后观察不再成立，就不能继续套原来的写法。

\textbf{相邻状态。} 规律是冷冻期把普通买卖股票拆成多个日末状态，转移不能跨过限制。把这个特例继续拆开看：先问暴力搜索里哪些不同路径到达同一个后续问题，再给这个后续问题命名为状态。状态必须包含未来决策需要的全部信息，转移只从更小且已经算清的状态来。

\textbf{状态复用。} 把重复递归节点命名为状态；遍历顺序只是在保证依赖状态已经算完。本题落点是：规律是冷冻期把普通买卖股票拆成多个日末状态，转移不能跨过限制。手算时只改被新输入影响的那一小部分，而不是回头重算完整候选。

\textbf{指针和字段。} 状态字段要能说清：状态下标、状态值语义、初始化、转移来源、遍历顺序；空间压缩时还要指出读到的是旧值还是新值。手算时按这个协议跟踪：拿 [1, 2, 3, 0, 2] 手算，第 2 天卖出后第 3 天不能立刻买入，因为有冷冻期；状态必须区分持股、刚卖出、空仓可买。然后按协议复查：手算时先写一个很小的 DP 表或记忆化搜索记录。每个格子旁边写它来自哪些更小状态。

\textbf{代码落点。} 先写未压缩版本校验转移；压缩前后用小表对拍；不可达哨兵参与加法前要检查溢出。关键代码行必须能逐句翻译成“旧状态怎样变成新状态”，否则就只是模板。

\textbf{正确性。} 证明从最后一步开始：任意最优解的最后一步都在转移枚举中，转移来源已经由更小状态正确计算。

\textbf{边界实验。} 先用空价格、一天、连续上涨、卖出后不能次日买入做反例检查。实验时覆盖空价格、一天、连续上涨、卖出后不能次日买入，先打印完整 DP 表，再比较空间压缩版本。若压缩版不同，优先检查遍历顺序和旧值保存。

\textbf{复杂度变化。} 暴力递归会重复计算相同子问题，常见指数级；DP 把每个状态算一次，时间是状态数乘单个转移成本，空间是状态表大小或压缩后的大小。

\textbf{迁移追问。} 如果题目改成“手续费版本可以在买入或卖出转移中扣费”，先判断原状态是否仍然覆盖新答案集合，再决定是微调边界、改 value 语义，还是换算法。

\begin{principlebox}{本节复盘方式}
不要只看最终算法名。每道题复盘时先遮住优化版，口述暴力枚举对象；再指出相邻窗口、历史前缀、候选区间、搜索状态或子问题里哪一部分被重复处理；最后用一个边界样例验证状态更新顺序。能讲完这三步，才算真正掌握本章案例。
\end{principlebox}
